% Lean module: MonadFormal.Temporal.MixedBoundary
\subsection{A literal mixed equation and a resonance obstruction}
\label{sec:mixed-boundary}

The eliminated-state identities do not settle well-posedness of an
arbitrary equation containing advanced arguments. Consider the scalar
boundary problem
\begin{equation}\label{eq:mixed-boundary}
 x'(t)=\alpha\bigl(x(t-\tfrac12)+x(t+\tfrac12)\bigr),\qquad
 x(t)=0\ (t\le0),\quad x(t)=1\ (t>1).
\end{equation}
Here a solution means a real-valued path continuous on $[0,1]$, with
the displayed derivative on the two open intervals $(0,\tfrac12)$ and
$(\tfrac12,1)$. We do not require differentiability at $\tfrac12$ or
continuity across the future boundary at $1$. These qualifications are
part of the problem, not conclusions of the calculation.
\leandef{MonadFormal.Temporal.MixedBoundary.IsSolution}

Put $u(s)=x(s)$ and $v(s)=x(s+\tfrac12)$ for $0\le s\le\tfrac12$.
The two open-interval equations become
\[
 u'=\alpha v,\qquad v'=\alpha(u+1),\qquad
 u(0)=0,\quad u(\tfrac12)=v(0).
\]
The last equality is the seam condition: both sides describe
$x(\tfrac12)$. Define
\begin{align*}
 u_\beta(s)&=\cosh(\alpha s)-1+\beta\sinh(\alpha s),\\
 v_\beta(s)&=\sinh(\alpha s)+\beta\cosh(\alpha s).
\end{align*}
\leandef{MonadFormal.Temporal.MixedBoundary.left,MonadFormal.Temporal.MixedBoundary.right}

\begin{lemma}[Exhaustive split-system solution]
\label[lemma]{lem:mixed-split}
Suppose $u,v$ are continuous on $[0,t]$, differentiable on $(0,t)$,
satisfy $u'=\alpha v$, $v'=\alpha(u+1)$ there, and $u(0)=0$.
For $t\ge0$, their endpoint values are
$u(t)=u_\beta(t)$ and $v(t)=v_\beta(t)$, where $\beta=v(0)$.
Conversely, $u_\beta,v_\beta$ satisfy the two differential equations
everywhere and have initial values $0,\beta$.
\end{lemma}
\begin{proof}
For $z'=cz$, the derivative of $z(s)e^{-cs}$ is zero on the
interior. Continuity at the endpoints and the fundamental theorem of
calculus give $z(t)=z(0)e^{ct}$. Apply this first to $u+1+v$, with
$c=\alpha$, and then to $u+1-v$, with $c=-\alpha$. Adding and
subtracting the resulting identities gives the displayed hyperbolic
formulas. Their derivatives and initial values follow directly from
the derivatives of $\sinh$ and $\cosh$.
\leanproof{MonadFormal.Temporal.MixedBoundary.scalar_solution,MonadFormal.Temporal.MixedBoundary.split_solution,MonadFormal.Temporal.MixedBoundary.left_zero,MonadFormal.Temporal.MixedBoundary.right_zero,MonadFormal.Temporal.MixedBoundary.left_derivative,MonadFormal.Temporal.MixedBoundary.right_derivative}
\end{proof}

\begin{theorem}[Boundary compatibility and exact existence criterion]
\label{thm:mixed-existence}
Write $S=\sinh(\alpha/2)$ and $C=\cosh(\alpha/2)$.
The seam condition is equivalent to
\[
 (1-S)\beta=C-1.
\]
If $S\ne1$, it selects exactly one coefficient,
$\beta=(C-1)/(1-S)$, and \eqref{eq:mixed-boundary} has a solution in
the declared class. If $S=1$, there is no solution in that class.
In particular, $\alpha=2\operatorname{arsinh}(1)$ is an explicit
resonant parameter at which existence fails.
\end{theorem}
\begin{proof}
Substitute the explicit formulas into $u_\beta(\tfrac12)=v_\beta(0)$.
When $S\ne1$, divide by $1-S$ and assemble
\[
 x_\beta(t)=
 \begin{cases}
 0,&t\le0,\\
 u_\beta(t),&0<t\le\tfrac12,\\
 v_\beta(t-\tfrac12),&\tfrac12<t\le1,\\
 1,&t>1.
 \end{cases}
\]
The initial and seam equalities make this path continuous on $[0,1]$.
On the first open half its delayed value is zero and its advanced
value is $v_\beta(t)$; on the second, those values are
$u_\beta(t-\tfrac12)$ and one. The split derivatives therefore give
\eqref{eq:mixed-boundary} on both open halves.

At $S=1$, the seam equation would require $C=1$, contradicting
$C^2-S^2=1$. This excludes every solution, not just the explicit
construction: any solution gives continuous split paths, and
\Cref{lem:mixed-split} forces their values at the seam to have the
same form. Finally, $\sinh(\operatorname{arsinh}(1))=1$ supplies the
stated parameter.
\leandef{MonadFormal.Temporal.MixedBoundary.path}
\leanproof{MonadFormal.Temporal.MixedBoundary.join_iff,MonadFormal.Temporal.MixedBoundary.join_unique,MonadFormal.Temporal.MixedBoundary.no_join_at_resonance,MonadFormal.Temporal.MixedBoundary.resonant_parameter,MonadFormal.Temporal.MixedBoundary.no_split_at_resonance,MonadFormal.Temporal.MixedBoundary.path_history,MonadFormal.Temporal.MixedBoundary.path_left,MonadFormal.Temporal.MixedBoundary.path_right,MonadFormal.Temporal.MixedBoundary.path_future,MonadFormal.Temporal.MixedBoundary.path_equation_left,MonadFormal.Temporal.MixedBoundary.path_equation_right,MonadFormal.Temporal.MixedBoundary.path_continuousOn,MonadFormal.Temporal.MixedBoundary.exists_mixed_solution,MonadFormal.Temporal.MixedBoundary.no_mixed_at_resonance,MonadFormal.Temporal.MixedBoundary.solution_exists_iff}
\end{proof}

\begin{corollary}[Matching the future endpoint is an extra constraint]
For the joined, non-resonant construction, $v_\beta(\tfrac12)=1$
if and only if $C=2S$, equivalently $\alpha=\log3$.
\end{corollary}
\begin{proof}
Multiply $v_\beta(\tfrac12)=S+\beta C$ by $1-S$, use the seam
equation and $C^2-S^2=1$, and obtain
\[
 v_\beta(\tfrac12)(1-S)=1+S-C.
\]
Since $1-S\ne0$, the desired value is equivalent to $C=2S$.
Writing the hyperbolic functions as exponentials then gives
$e^\alpha=3$, equivalently $\alpha=\log3$.
\leanproof{MonadFormal.Temporal.MixedBoundary.terminal_match_iff,MonadFormal.Temporal.MixedBoundary.terminal_parameter_iff}
\end{proof}

Thus the solution convention matters. A path in the declared class may
jump immediately after $1$; requiring a continuous join to the unit
future imposes the additional endpoint equation. Neither the explicit
construction nor its resonance obstruction establishes a general
well-posedness theorem for advanced functional differential equations.
